\appendix

\section{Supplementary Empirical Accounting}
\label{app:supplementary}

The appendices report supporting accounting and diagnostic views of the frozen
empirical evidence. They do not introduce additional repository executions,
mutation outcomes, or post hoc case selection.

\subsection{Full frozen corpus}
\label{app:full-corpus}

Table~\ref{tab:full-corpus} lists all 39 frozen repository--operator cases used
in the empirical study. The table joins the final accounting layer with the
outcome-blind workflow frame by frozen case identifier. B01 oracle entries are
left unclassified because oracle kind was not prospectively recorded under the
schema used for that calibration batch.

The displayed outcome is the combined study result: the canonical primary
result when mutation evaluation succeeded at the primary layer, the bounded
D027 result when a previously non-evaluable frozen case was successfully
restored and evaluated, and non-evaluable otherwise. This presentation does
not alter the separate primary and restoration accounting reported in the main
text.

\begingroup
\scriptsize
\setlength{\tabcolsep}{2pt}
\renewcommand{\arraystretch}{1.15}
\begin{longtable}{@{}>{\raggedright\arraybackslash}p{0.075\textwidth}>{\raggedright\arraybackslash}p{0.245\textwidth}>{\raggedright\arraybackslash}p{0.115\textwidth}>{\raggedright\arraybackslash}p{0.145\textwidth}>{\raggedright\arraybackslash}p{0.110\textwidth}>{\raggedright\arraybackslash}p{0.070\textwidth}>{\raggedright\arraybackslash}p{0.110\textwidth}@{}}
\caption{Full frozen 39-case empirical corpus. Oracle classification is shown only where it was prospectively recorded for B02; B01 entries are left unclassified. Source identifies whether the combined result comes from the canonical primary execution or the bounded D027 restoration layer.}\label{tab:full-corpus}\\
\toprule
\textbf{Case} &
\textbf{Repository} &
\textbf{Operator} &
\textbf{Workflow} &
\textbf{Oracle} &
\textbf{Source} &
\textbf{Outcome} \\
\midrule
\endfirsthead

\multicolumn{7}{c}{\tablename\ \thetable\ --- continued from previous page} \\
\toprule
\textbf{Case} &
\textbf{Repository} &
\textbf{Operator} &
\textbf{Workflow} &
\textbf{Oracle} &
\textbf{Source} &
\textbf{Outcome} \\
\midrule
\endhead

\midrule
\multicolumn{7}{r}{Continued on next page} \\
\endfoot

\bottomrule
\endlastfoot

\texttt{B01-01} & \texttt{snu}\texttt{-}\allowbreak{}\texttt{causality}\texttt{-}\allowbreak{}\texttt{lab}\texttt{/}\allowbreak{}\texttt{efficient}\texttt{-}\allowbreak{}\texttt{canonical}\texttt{-}\allowbreak{}\texttt{bounding} & \texttt{dependency}\texttt{-}\allowbreak{}\texttt{pin} & \texttt{upstream}\texttt{-}\allowbreak{}\texttt{test} & -- & primary & \textsc{Survived} \\
\texttt{B01-02} & \texttt{scikit}\texttt{-}\allowbreak{}\texttt{learn}\texttt{-}\allowbreak{}\texttt{contrib}\texttt{/}\allowbreak{}\texttt{imbalanced}\texttt{-}\allowbreak{}\texttt{learn} & \texttt{data}\texttt{-}\allowbreak{}\texttt{split} & \texttt{documented}\texttt{-}\allowbreak{}\texttt{example} & -- & D027 & \textsc{Survived} \\
\texttt{B01-03} & \texttt{scikit}\texttt{-}\allowbreak{}\texttt{learn}\texttt{-}\allowbreak{}\texttt{contrib}\texttt{/}\allowbreak{}\texttt{MAPIE} & \texttt{random}\texttt{-}\allowbreak{}\texttt{seed} & \texttt{documented}\texttt{-}\allowbreak{}\texttt{example} & -- & primary & \textsc{Survived} \\
\texttt{B01-04} & \texttt{wwu}\texttt{-}\allowbreak{}\texttt{mmll}\texttt{/}\allowbreak{}\texttt{photonai} & \texttt{cv}\texttt{-}\allowbreak{}\texttt{fold}\texttt{-}\allowbreak{}\texttt{count} & \texttt{documented}\texttt{-}\allowbreak{}\texttt{example} & -- & primary & \textsc{Survived} \\
\texttt{B01-05} & \texttt{damianhorna}\texttt{/}\allowbreak{}\texttt{multi}\texttt{-}\allowbreak{}\texttt{imbalance} & \texttt{data}\texttt{-}\allowbreak{}\texttt{split} & \texttt{upstream}\texttt{-}\allowbreak{}\texttt{test} & -- & primary & \textsc{Survived} \\
\texttt{B01-06} & \texttt{KoheiObata}\texttt{/}\allowbreak{}\texttt{DMM} & \texttt{dependency}\texttt{-}\allowbreak{}\texttt{pin} & \texttt{documented}\texttt{-}\allowbreak{}\texttt{experiment} & -- & -- & non-\allowbreak{}evaluable \\
\texttt{B01-07} & \texttt{sigeisler}\texttt{/}\allowbreak{}\texttt{robustness}\texttt{\_}\allowbreak{}\texttt{of}\texttt{\_}\allowbreak{}\texttt{gnns}\texttt{\_}\allowbreak{}\texttt{at}\texttt{\_}\allowbreak{}\texttt{scale} & \texttt{dependency}\texttt{-}\allowbreak{}\texttt{pin} & \texttt{upstream}\texttt{-}\allowbreak{}\texttt{test} & -- & -- & non-\allowbreak{}evaluable \\
\texttt{B01-08} & \texttt{ContextLab}\texttt{/}\allowbreak{}\texttt{hypertools} & \texttt{random}\texttt{-}\allowbreak{}\texttt{seed} & \texttt{documented}\texttt{-}\allowbreak{}\texttt{example} & -- & primary & \textsc{Survived} \\
\texttt{B01-09} & \texttt{sherbold}\texttt{/}\allowbreak{}\texttt{autorank} & \texttt{random}\texttt{-}\allowbreak{}\texttt{seed} & \texttt{documented}\texttt{-}\allowbreak{}\texttt{example} & -- & primary & \textsc{Survived} \\
\texttt{B01-10} & \texttt{BorgwardtLab}\texttt{/}\allowbreak{}\texttt{P}\texttt{-}\allowbreak{}\texttt{WL} & \texttt{cv}\texttt{-}\allowbreak{}\texttt{fold}\texttt{-}\allowbreak{}\texttt{count} & \texttt{documented}\texttt{-}\allowbreak{}\texttt{experiment} & -- & primary & \textsc{Survived} \\
\texttt{B02-01} & \texttt{tslearn}\texttt{-}\allowbreak{}\texttt{team}\texttt{/}\allowbreak{}\texttt{tslearn} & \texttt{random}\texttt{-}\allowbreak{}\texttt{seed} & \texttt{documented}\texttt{-}\allowbreak{}\texttt{example} & \texttt{completion}\texttt{-}\allowbreak{}\texttt{only} & primary & \textsc{Survived} \\
\texttt{B02-02} & \texttt{tensorly}\texttt{/}\allowbreak{}\texttt{tensorly} & \texttt{random}\texttt{-}\allowbreak{}\texttt{seed} & \texttt{documented}\texttt{-}\allowbreak{}\texttt{example} & \texttt{completion}\texttt{-}\allowbreak{}\texttt{only} & primary & \textsc{Survived} \\
\texttt{B02-03} & \texttt{braindecode}\texttt{/}\allowbreak{}\texttt{braindecode} & \texttt{random}\texttt{-}\allowbreak{}\texttt{seed} & \texttt{documented}\texttt{-}\allowbreak{}\texttt{example} & \texttt{completion}\texttt{-}\allowbreak{}\texttt{only} & D027 & \textsc{Survived} \\
\texttt{B02-04} & \texttt{rtqichen}\texttt{/}\allowbreak{}\texttt{torchdiffeq} & \texttt{random}\texttt{-}\allowbreak{}\texttt{seed} & \texttt{documented}\texttt{-}\allowbreak{}\texttt{example} & \texttt{completion}\texttt{-}\allowbreak{}\texttt{only} & primary & \textsc{Survived} \\
\texttt{B02-05} & \texttt{EthanJamesLew}\texttt{/}\allowbreak{}\texttt{AutoKoopman} & \texttt{random}\texttt{-}\allowbreak{}\texttt{seed} & \texttt{documented}\texttt{-}\allowbreak{}\texttt{experiment} & \texttt{completion}\texttt{-}\allowbreak{}\texttt{only} & -- & non-\allowbreak{}evaluable \\
\texttt{B02-06} & \texttt{DistrictDataLabs}\texttt{/}\allowbreak{}\texttt{yellowbrick} & \texttt{random}\texttt{-}\allowbreak{}\texttt{seed} & \texttt{upstream}\texttt{-}\allowbreak{}\texttt{test} & \texttt{assertion} & -- & non-\allowbreak{}evaluable \\
\texttt{B02-07} & \texttt{kLabUM}\texttt{/}\allowbreak{}\texttt{rrcf} & \texttt{random}\texttt{-}\allowbreak{}\texttt{seed} & \texttt{upstream}\texttt{-}\allowbreak{}\texttt{test} & \texttt{assertion} & D027 & \textsc{Survived} \\
\texttt{B02-08} & \texttt{eric}\texttt{-}\allowbreak{}\texttt{mitchell}\texttt{/}\allowbreak{}\texttt{detect}\texttt{-}\allowbreak{}\texttt{gpt} & \texttt{random}\texttt{-}\allowbreak{}\texttt{seed} & \texttt{documented}\texttt{-}\allowbreak{}\texttt{experiment} & \texttt{completion}\texttt{-}\allowbreak{}\texttt{only} & -- & non-\allowbreak{}evaluable \\
\texttt{B02-09} & \texttt{gnina}\texttt{/}\allowbreak{}\texttt{libmolgrid} & \texttt{random}\texttt{-}\allowbreak{}\texttt{seed} & \texttt{upstream}\texttt{-}\allowbreak{}\texttt{test} & \texttt{assertion} & -- & non-\allowbreak{}evaluable \\
\texttt{B02-10} & \texttt{xucong}\texttt{-}\allowbreak{}\texttt{zhang}\texttt{/}\allowbreak{}\texttt{ETH}\texttt{-}\allowbreak{}\texttt{XGaze} & \texttt{random}\texttt{-}\allowbreak{}\texttt{seed} & \texttt{documented}\texttt{-}\allowbreak{}\texttt{experiment} & \texttt{completion}\texttt{-}\allowbreak{}\texttt{only} & -- & non-\allowbreak{}evaluable \\
\texttt{B02-11} & \texttt{princeton}\texttt{-}\allowbreak{}\texttt{nlp}\texttt{/}\allowbreak{}\texttt{SimCSE} & \texttt{dependency}\texttt{-}\allowbreak{}\texttt{pin} & \texttt{documented}\texttt{-}\allowbreak{}\texttt{validation} & \texttt{completion}\texttt{-}\allowbreak{}\texttt{only} & -- & non-\allowbreak{}evaluable \\
\texttt{B02-12} & \texttt{sunblaze}\texttt{-}\allowbreak{}\texttt{ucb}\texttt{/}\allowbreak{}\texttt{dpml}\texttt{-}\allowbreak{}\texttt{benchmark} & \texttt{dependency}\texttt{-}\allowbreak{}\texttt{pin} & \texttt{documented}\texttt{-}\allowbreak{}\texttt{experiment} & \texttt{completion}\texttt{-}\allowbreak{}\texttt{only} & -- & non-\allowbreak{}evaluable \\
\texttt{B02-13} & \texttt{rcamino}\texttt{/}\allowbreak{}\texttt{multi}\texttt{-}\allowbreak{}\texttt{categorical}\texttt{-}\allowbreak{}\texttt{gans} & \texttt{dependency}\texttt{-}\allowbreak{}\texttt{pin} & \texttt{documented}\texttt{-}\allowbreak{}\texttt{experiment} & \texttt{completion}\texttt{-}\allowbreak{}\texttt{only} & -- & non-\allowbreak{}evaluable \\
\texttt{B02-14} & \texttt{IBM}\texttt{/}\allowbreak{}\texttt{concept}\texttt{\_}\allowbreak{}\texttt{transformer} & \texttt{dependency}\texttt{-}\allowbreak{}\texttt{pin} & \texttt{upstream}\texttt{-}\allowbreak{}\texttt{test} & \texttt{assertion} & primary & \textsc{Equivalent} \\
\texttt{B02-15} & \texttt{apple}\texttt{/}\allowbreak{}\texttt{ml}\texttt{-}\allowbreak{}\texttt{cvpr2019}\texttt{-}\allowbreak{}\texttt{swd} & \texttt{dependency}\texttt{-}\allowbreak{}\texttt{pin} & \texttt{documented}\texttt{-}\allowbreak{}\texttt{example} & \texttt{completion}\texttt{-}\allowbreak{}\texttt{only} & D027 & \textsc{Killed} \\
\texttt{B02-16} & \texttt{GestureGeneration}\texttt{/}\allowbreak{}\texttt{Speech}\texttt{\_}\allowbreak{}\texttt{driven}\texttt{\_}\allowbreak{}\texttt{gesture}\texttt{\_}\allowbreak{}\texttt{generation}\texttt{\_}\allowbreak{}\texttt{with}\texttt{\_}\allowbreak{}\texttt{autoencoder} & \texttt{dependency}\texttt{-}\allowbreak{}\texttt{pin} & \texttt{documented}\texttt{-}\allowbreak{}\texttt{experiment} & \texttt{completion}\texttt{-}\allowbreak{}\texttt{only} & -- & non-\allowbreak{}evaluable \\
\texttt{B02-17} & \texttt{AaltoML}\texttt{/}\allowbreak{}\texttt{BayesNewton} & \texttt{dependency}\texttt{-}\allowbreak{}\texttt{pin} & \texttt{upstream}\texttt{-}\allowbreak{}\texttt{test} & \texttt{assertion} & -- & non-\allowbreak{}evaluable \\
\texttt{B02-18} & \texttt{pulimeng}\texttt{/}\allowbreak{}\texttt{eToxPred} & \texttt{dependency}\texttt{-}\allowbreak{}\texttt{pin} & \texttt{documented}\texttt{-}\allowbreak{}\texttt{validation} & \texttt{completion}\texttt{-}\allowbreak{}\texttt{only} & D027 & \textsc{Survived} \\
\texttt{B02-19} & \texttt{ddd9898}\texttt{/}\allowbreak{}\texttt{DeepNano} & \texttt{dependency}\texttt{-}\allowbreak{}\texttt{pin} & \texttt{documented}\texttt{-}\allowbreak{}\texttt{validation} & \texttt{completion}\texttt{-}\allowbreak{}\texttt{only} & primary & \textsc{Killed} \\
\texttt{B02-20} & \texttt{MKMaS}\texttt{-}\allowbreak{}\texttt{GUET}\texttt{/}\allowbreak{}\texttt{SSJE} & \texttt{dependency}\texttt{-}\allowbreak{}\texttt{pin} & \texttt{documented}\texttt{-}\allowbreak{}\texttt{experiment} & \texttt{completion}\texttt{-}\allowbreak{}\texttt{only} & -- & non-\allowbreak{}evaluable \\
\texttt{B02-21} & \texttt{thomas}\texttt{-}\allowbreak{}\texttt{young}\texttt{-}\allowbreak{}\texttt{2013}\texttt{/}\allowbreak{}\texttt{mindware} & \texttt{data}\texttt{-}\allowbreak{}\texttt{split} & \texttt{ci} & \texttt{completion}\texttt{-}\allowbreak{}\texttt{only} & D027 & \textsc{Survived} \\
\texttt{B02-22} & \texttt{Alcoholrithm}\texttt{/}\allowbreak{}\texttt{PTaRL} & \texttt{data}\texttt{-}\allowbreak{}\texttt{split} & \texttt{upstream}\texttt{-}\allowbreak{}\texttt{test} & \texttt{completion}\texttt{-}\allowbreak{}\texttt{only} & D027 & \textsc{Survived} \\
\texttt{B02-23} & \texttt{tfmortie}\texttt{/}\allowbreak{}\texttt{setvaluedprediction} & \texttt{data}\texttt{-}\allowbreak{}\texttt{split} & \texttt{upstream}\texttt{-}\allowbreak{}\texttt{test} & \texttt{completion}\texttt{-}\allowbreak{}\texttt{only} & D027 & \textsc{Survived} \\
\texttt{B02-24} & \texttt{lucazav}\texttt{/}\allowbreak{}\texttt{binclass}\texttt{-}\allowbreak{}\texttt{tools} & \texttt{data}\texttt{-}\allowbreak{}\texttt{split} & \texttt{upstream}\texttt{-}\allowbreak{}\texttt{test} & \texttt{assertion} & -- & non-\allowbreak{}evaluable \\
\texttt{B02-25} & \texttt{ECOLE}\texttt{-}\allowbreak{}\texttt{ITN}\texttt{/}\allowbreak{}\texttt{NguyenSSCI2021} & \texttt{data}\texttt{-}\allowbreak{}\texttt{split} & \texttt{upstream}\texttt{-}\allowbreak{}\texttt{test} & \texttt{completion}\texttt{-}\allowbreak{}\texttt{only} & -- & non-\allowbreak{}evaluable \\
\texttt{B02-26} & \texttt{tungtokyo1108}\texttt{/}\allowbreak{}\texttt{PolicySynth} & \texttt{data}\texttt{-}\allowbreak{}\texttt{split} & \texttt{documented}\texttt{-}\allowbreak{}\texttt{example} & \texttt{completion}\texttt{-}\allowbreak{}\texttt{only} & primary & \textsc{Survived} \\
\texttt{B02-27} & \texttt{rsteca}\texttt{/}\allowbreak{}\texttt{sklearn}\texttt{-}\allowbreak{}\texttt{deap} & \texttt{cv}\texttt{-}\allowbreak{}\texttt{fold}\texttt{-}\allowbreak{}\texttt{count} & \texttt{upstream}\texttt{-}\allowbreak{}\texttt{test} & \texttt{assertion} & D027 & \textsc{Survived} \\
\texttt{B02-28} & \texttt{ogozuacik}\texttt{/}\allowbreak{}\texttt{d3}\texttt{-}\allowbreak{}\texttt{discriminative}\texttt{-}\allowbreak{}\texttt{drift}\texttt{-}\allowbreak{}\texttt{detector}\texttt{-}\allowbreak{}\texttt{concept}\texttt{-}\allowbreak{}\texttt{drift} & \texttt{cv}\texttt{-}\allowbreak{}\texttt{fold}\texttt{-}\allowbreak{}\texttt{count} & \texttt{documented}\texttt{-}\allowbreak{}\texttt{experiment} & \texttt{completion}\texttt{-}\allowbreak{}\texttt{only} & D027 & \textsc{Survived} \\
\texttt{B02-29} & \texttt{farshidrayhancv}\texttt{/}\allowbreak{}\texttt{CUSBoost} & \texttt{cv}\texttt{-}\allowbreak{}\texttt{fold}\texttt{-}\allowbreak{}\texttt{count} & \texttt{documented}\texttt{-}\allowbreak{}\texttt{experiment} & \texttt{completion}\texttt{-}\allowbreak{}\texttt{only} & D027 & \textsc{Survived} \\
\end{longtable}
\endgroup


\subsection{Bounded restoration accounting}
\label{app:d027}

The primary empirical layer remains the canonical execution record. D027 was a
separate bounded restoration layer intended to recover evaluability for frozen
cases without changing the selected repository revision, mutation candidate,
validation workflow, or oracle.

Within the 24-case D027 cohort, 19 cases had a restoration report or assessment.
Substantive restoration was attempted for 18 cases. Eleven cases were restored
and reached mutation evaluation, seven were not restored after a substantive
attempt, and one additional reported case was not substantively attempted
because repository identity could not be recovered.

Table~\ref{tab:d027-accounting} reports this accounting separately from the
primary mutation outcomes.

\begin{table}[t]
\centering
\caption{D027 bounded-restoration accounting. Report presence and substantive restoration attempt are distinct.}
\label{tab:d027-accounting}
\small
\renewcommand{\arraystretch}{1.18}
\setlength{\tabcolsep}{7pt}
\begin{tabular}{lrr}
\toprule
\textbf{State} & \textbf{Count} & \textbf{\% of cohort} \\
\midrule
Cohort & 24 & 100.0\% \\
Report/assessment present & 19 & 79.2\% \\
No D027 report & 5 & 20.8\% \\
Substantive restoration attempted & 18 & 75.0\% \\
Report present, restoration not attempted & 1 & 4.2\% \\
Restored & 11 & 45.8\% \\
Not restored after substantive attempt & 7 & 29.2\% \\
Mutation evaluated & 11 & 45.8\% \\
\bottomrule
\end{tabular}
\end{table}



\subsection{Operator-by-oracle diagnostic}
\label{app:operator-oracle}

RQ2 was analyzed descriptively because the oracle categories were sparse and
operator composition was uneven. Table~\ref{tab:rq2-operator-oracle} reports
the B02 operator-by-oracle cross-classification used as a diagnostic for this
confounding.

In particular, both detected B02 mutations occurred in
\texttt{dependency-pin} cases classified as \texttt{completion-only}. The table
is therefore useful for showing why the observed oracle-level proportions
cannot be interpreted independently of mutation operator.

\begin{table}[t]
\centering
\caption{Operator composition of the prospective B02 oracle analysis. Both detected mutations occurred in the dependency-pin/completion-only cell; sparse cells prevent attribution of this pattern to oracle kind.}
\label{tab:rq2-operator-oracle}
\begin{tabular}{lrrrr}
\toprule
Operator & Oracle kind & N & Killed & Survived \\
\midrule
random-seed & assertion & 1 & 0 & 1 \\
random-seed & completion-only & 4 & 0 & 4 \\
dependency-pin & completion-only & 3 & 2 & 1 \\
data-split & completion-only & 4 & 0 & 4 \\
cv-fold-count & assertion & 1 & 0 & 1 \\
cv-fold-count & completion-only & 2 & 0 & 2 \\
\bottomrule
\end{tabular}
\end{table}



\subsection{Evaluability by workflow and oracle category}
\label{app:rq2-attrition}

Mutation outcomes are observed only for cases that reach successful baseline
execution and mutation evaluation. Differential evaluability can therefore
change the composition of the analyzed subsets.

Table~\ref{tab:rq2-oracle-attrition} reports selected, evaluated, and meaningful case
counts by the categories used in RQ2. These counts are presented to make
attrition visible rather than implicitly conditioning the analysis on only
successful executions.

\begin{table}[t]
\centering
\caption{Evaluability of prospectively classified B02 oracle categories.}
\label{tab:rq2-oracle-attrition}
\begin{tabular}{lrrr}
\toprule
Oracle kind & Selected & Combined evaluated & Meaningful \\
\midrule
assertion & 7 & 3 & 2 \\
completion-only & 22 & 13 & 13 \\
\bottomrule
\end{tabular}
\end{table}



\section{Mutation Operator Specification}
\label{app:operator-specification}

The empirical study used the concise labels
\texttt{random-seed},
\texttt{dependency-pin},
\texttt{data-split}, and
\texttt{cv-fold-count}. The corresponding implemented transformations were
deliberately narrow and deterministic once a mutation candidate had been
selected.

\begin{description}

\item[\texttt{random-seed}.]
Targets supported integer-literal seed calls and replaces the selected seed
\(N\) with \(N+1\). Supported calls in the study implementation include
Python's \texttt{random.seed}, NumPy seed calls, and
\texttt{torch.manual\_seed}.

\item[\texttt{dependency-pin}.]
Targets an exact dependency requirement of the form
\texttt{package==version} and replaces the exact constraint with
\texttt{package>=version}. Dependency evaluation independently resolves the
baseline and mutant environments. A textual mutation that resolves to the same
installed target version is treated as semantically equivalent.

\item[\texttt{data-split}.]
Targets a supported scikit-learn \texttt{train\_test\_split} call containing an
explicit non-\texttt{None} \texttt{stratify} argument and replaces the
stratification expression with \texttt{None}.

\item[\texttt{cv-fold-count}.]
Targets an explicit literal \texttt{n\_splits} value in a supported
KFold-family splitter and replaces \(N\) with \(N+1\). Supported splitters are
\texttt{KFold}, \texttt{StratifiedKFold}, \texttt{RepeatedKFold}, and
\texttt{RepeatedStratifiedKFold}.

\end{description}

Candidate detection and mutation application are separate operations.
Source-level candidates retain target-location metadata and are matched again
against the source before mutation application. This prevents a selected
candidate from being silently redirected to another syntactic occurrence.


\section{Outcome and Evidence Semantics}
\label{app:outcome-semantics}

The empirical analysis distinguishes execution state, mutation outcome, and
semantic verification.

A \textsc{Killed} mutation is one for which the frozen validation workflow
returns a failing status after a successful baseline. A \textsc{Survived}
mutation is one for which the same workflow succeeds after mutation.
Infrastructure, setup, baseline, and restoration failures are not mutation
outcomes.

Semantic verification is recorded independently. The relevant study-level
states are confirmed non-equivalent, confirmed equivalent, unverified, and not
run. Confirmed-equivalent mutations are excluded from the meaningful detection
denominator.

Consequently, the main detection measure is

\[
\frac{
  \textsc{Killed}
}{
  \textsc{Killed}
  +
  \text{confirmed non-equivalent \textsc{Survived}}
}.
\]

A surviving mutation therefore has a deliberately narrow interpretation: the
selected frozen validation workflow did not detect that particular confirmed
non-equivalent mutation. It is not a repository-level classification of
reproducibility.


\section{Empirical Provenance and Freeze Boundaries}
\label{app:freeze-boundaries}

The empirical study used explicit freeze boundaries to separate corpus
construction, execution, restoration, and analysis.

The 39-case empirical frame was frozen before post-freeze analysis. Candidate
repositories were not added, replaced, or rerun to improve the observed
mutation-detection rate. D027 restoration retained the frozen revision,
candidate, workflow, and oracle and did not rewrite the canonical primary
execution record.

For RQ2, workflow categories came from the frozen corpus metadata. Oracle
categories for the B02 cases were defined prospectively under protocol version
2 and frozen before joining them with mutation outcomes. The B01 calibration
cases predated this schema and were therefore not retrospectively assigned
oracle categories.

The repository contains the machine-readable case records, restoration
assessments, RQ2 blind frame, generated result tables, and manifests used to
produce the reported accounting. These artifacts are retained as the auditable
record underlying the manuscript tables.
